\section{Introduction}\label{sec:introduction}

\begin{figure*}[!t]
  \centering
  \resizebox{\linewidth}{!}{\begin{tikzpicture}[petriBase, inline, scale=0.95, every node/.style={transform shape},
    place/.append style={minimum size=6mm},
    transition/.append style={minimum size=6.5mm, font=\small},
    plab/.style={font=\footnotesize, black!60, inner sep=1pt},
    ptitle/.style={font=\small\bfseries, inner sep=2pt},
    note/.style={font=\scriptsize, black!55, inner sep=1pt},
    rule/.style={black!30, line width=0.5pt},
    cnact/.style={rectangle, rounded corners=2pt, draw=black!55, fill=black!5,
                  minimum size=6.5mm, font=\small},
    omk/.style={circle, draw=#1, fill=#1, minimum size=4.2pt, inner sep=0pt},
    smk/.style={rectangle, draw=#1, fill=#1, minimum size=4.8pt, inner sep=0pt},
    card/.style={font=\scriptsize, inner sep=0.5pt},
    bind/.style={black!55, line width=0.5pt, dashed},
    op/.style={circle, draw=black!55, fill=black!4, inner sep=1.2pt, minimum size=5mm,
               font=\small},
    lboth/.style={rectangle, rounded corners=2pt, draw=black!55, fill=black!5,
                  minimum size=6.5mm, font=\small},
    litem/.style={rectangle, rounded corners=2pt, draw=redcol, fill=redcol!12,
                  minimum size=6.5mm, font=\small},
    task/.style={rectangle, draw=black!65, fill=black!4, minimum size=6.5mm, font=\small},
    gw/.style={diamond, draw=black!60, fill=black!6, inner sep=1pt, minimum size=7.6mm,
               font=\small},
    wt/.style={font=\scriptsize, inner sep=1pt},
    trace/.style={font=\small, inner sep=1.5pt},
    uwire/.style={typedwire, draw=black!75},
    sbox/.style={rectangle, draw=boxblueborder, thick, fill=boxblue, outer sep=0pt,
                 minimum width=0.44cm, minimum height=0.44cm, font=\scriptsize\sffamily}]

  \def\ycap{-1.95}
  \def\ycapB{-2.60}
  \def\rowB{-5.10}

  \providecommand{\mtfbrief}{0}
  \ifnum\mtfbrief=1
    \def\xSD{7.10}\def\ySD{0}
    \def\ycap{-2.30}
    \def\capSD{\ycap}
    \def\legendY{-3.55}
    \def\numOCPN{}\def\numSD{}\def\sufOCPN{}
  \else
    \def\xSD{12.35}\def\ySD{-5.10}
    \def\capSD{\ycapB}
    \def\numOCPN{(i) }\def\numSD{(vi) }\def\sufOCPN{ 1}
  \fi

  \begin{scope}[shift={(0,0)}]
    \node[T]                (a)   at (0,0)          {$a$};
    \node[blueP, tokens=1]  (po)  at (2.3,0.95)     {};
    \node[redP,  tokens=2]  (pi)  at (1.15,-0.95)   {};
    \node[T]                (p)   at (2.3,-0.95)    {$p$};
    \node[redP]             (pi2) at (3.45,-0.95)   {};
    \node[T]                (s)   at (4.6,0)        {$s$};

    \draw[blueflow] (a)   -- (po);
    \draw[blueflow] (po)  -- (s);
    \draw[redflow]  (a)   -- (pi) node[wt, pos=0.5, below left]  {$2$};
    \draw[redflow]  (pi)  -- (p);
    \draw[redflow]  (p)   -- (pi2);
    \draw[redflow]  (pi2) -- (s)  node[wt, pos=0.5, below right] {$2$};

    \ifnum\mtfbrief=1
      \node[note, anchor=north] at (2.3,-1.45) {arc weight};
    \fi
    \node[ptitle] at (2.3,\ycap) {\numOCPN object-centric Petri net\sufOCPN};
  \end{scope}

  \ifnum\mtfbrief=0
  \begin{scope}[shift={(5.90,-0.36)}]
    \node[cnact] (ca) at (0,0)     {$a$};
    \node[cnact] (cp) at (2.2,1.0) {$p$};
    \node[cnact] (cs) at (4.4,0)   {$s$};

    \draw[redflow,  name path=rap] (ca) to[out=48,in=180] (cp.west);
    \draw[redflow,  name path=rps] (cp.east) to[out=0,in=132] (cs);
    \draw[blueflow, name path=blu] (ca) to[out=-32,in=212] (cs);

    \path[name path=vA]    (0.85,-1.0) -- (0.85,1.7);
    \path[name path=vS]    (3.55,-1.0) -- (3.55,1.7);
    \path[name path=vPin]  (1.45,-1.0) -- (1.45,1.7);
    \path[name path=vPout] (2.95,-1.0) -- (2.95,1.7);
    \path[name intersections={of=rap and vA,    by=mItemA}];
    \path[name intersections={of=blu and vA,    by=mOrdA}];
    \path[name intersections={of=rps and vS,    by=mItemS}];
    \path[name intersections={of=blu and vS,    by=mOrdS}];
    \path[name intersections={of=rap and vPin,  by=pIn}];
    \path[name intersections={of=rps and vPout, by=pOut}];

    \draw[bind] (mItemA) -- (mOrdA);
    \draw[bind] (mItemS) -- (mOrdS);

    \node[smk=redcol, label={[card,redcol]above left:$2$}]  at (mItemA) {};
    \node[smk=redcol, label={[card,redcol]above right:$2$}] at (mItemS) {};
    \node[omk=redcol]  at (pIn)   {};
    \node[omk=redcol]  at (pOut)  {};
    \node[omk=bluecol] at (mOrdA) {};
    \node[omk=bluecol] at (mOrdS) {};

    \node[ptitle] at (2.2,{\ycap+0.36}) {(ii) object-centric causal net};
  \end{scope}
  \fi

  \ifnum\mtfbrief=0
  \begin{scope}[shift={(13.90,-0.39)}]
    \def\dd{0.06}
    \providecommand{\dualleg}[2]{%
      \draw[bluewire] ($(#1)!\dd cm!90:(#2)$)  -- ($(#2)!\dd cm!-90:(#1)$);
      \draw[redwire]  ($(#1)!\dd cm!-90:(#2)$) -- ($(#2)!\dd cm!90:(#1)$);
    }
    \coordinate (tr)   at (0,1.55);
    \coordinate (tpl)  at (0,0.45);
    \coordinate (tAt)  at (-1.37,0.67);
    \coordinate (tSt)  at (1.37,0.67);
    \coordinate (tP1t) at (-0.6,-0.48);
    \coordinate (tP2t) at (0.6,-0.48);
    \coordinate (trA)  at ($(tr)!0.20cm!(tAt)$);
    \coordinate (trS)  at ($(tr)!0.20cm!(tSt)$);

    \dualleg{tAt}{trA}
    \dualleg{trS}{tSt}
    \draw[redwire] (0,1.37) -- (0,0.62);
    \draw[redwire] ($(tpl)!0.20cm!(tP1t)$) -- (tP1t);
    \draw[redwire] ($(tpl)!0.20cm!(tP2t)$) -- (tP2t);

    \node[op]    (root)  at (tr)          {$\to$};
    \node[lboth] (ta)    at (-1.45,0.45)  {$a$};
    \node[op]    (tplus) at (tpl)         {$+$};
    \node[lboth] (ts)    at (1.45,0.45)   {$s$};
    \node[litem] (tp1)   at (-0.6,-0.70)  {$p$};
    \node[litem] (tp2)   at (0.6,-0.70)   {$p$};

    \node[ptitle] at (0,{\ycap+0.39}) {(iii) object-centric process tree};
  \end{scope}
  \fi

  \ifnum\mtfbrief=0
  \begin{scope}[shift={(5.80,\rowB+0.69)}]
    \node[task] (ba)  at (0,0)        {$a$};
    \node[gw]   (bg1) at (1.15,-0.90) {$+$};
    \node[task] (bp1) at (2.30,-0.20) {$p$};
    \node[task] (bp2) at (2.30,-1.60) {$p$};
    \node[gw]   (bg2) at (3.45,-0.90) {$+$};
    \node[task] (bs)  at (4.60,0)     {$s$};

    \draw[blueflow] (ba) to[out=25,in=155] (bs);
    \draw[redflow]  (ba)  -- (bg1);
    \draw[redflow]  (bg1) -- (bp1);
    \draw[redflow]  (bg1) -- (bp2);
    \draw[redflow]  (bp1) -- (bg2);
    \draw[redflow]  (bp2) -- (bg2);
    \draw[redflow]  (bg2) -- (bs);

    \node[ptitle] at (2.30,{\ycapB-0.69}) {(v) typed BPMN};
  \end{scope}

  \fi

  \begin{scope}[shift={(\xSD,\ySD)}]
    \node[morphismblue, minimum height=2cm, minimum width=0.95cm]    (da)  at (0,0)      {$a$};
    \node[morphismblue, minimum height=0.65cm, minimum width=0.65cm] (dp1) at (1.5,0.75) {$p$};
    \node[morphismblue, minimum height=0.65cm, minimum width=0.65cm] (dp2) at (1.5,0.0)  {$p$};
    \node[morphismblue, minimum height=2cm, minimum width=0.95cm]    (ds)  at (3.1,0)    {$s$};

    \draw[redwire] ($(da.south east)!0.875!(da.north east)$) to[out=0,in=180] (dp1.west);
    \draw[redwire] (dp1.east) to[out=0,in=180] ($(ds.south west)!0.875!(ds.north west)$);
    \draw[redwire] ($(da.south east)!0.5!(da.north east)$) to[out=0,in=180] (dp2.west);
    \draw[redwire] (dp2.east) to[out=0,in=180] ($(ds.south west)!0.5!(ds.north west)$);
    \draw[bluewire] ($(da.south east)!0.125!(da.north east)$) -- ($(ds.south west)!0.125!(ds.north west)$);

    \ifnum\mtfbrief=1
      \node[note, anchor=north] at (1.55,-1.45) {one wire per item};
    \fi
    \node[ptitle] at (1.55,\capSD) {\numSD string diagram};
  \end{scope}

  \ifnum\mtfbrief=0

  \begin{scope}[shift={(0,\rowB+0.18)}]
    \node[T]                (na)   at (0,0)          {$a$};
    \node[blueP, tokens=1]  (npo)  at (2.3,1.15)     {};
    \node[redP,  tokens=1]  (npi1) at (1.15,-0.30)   {};
    \node[T]                (np1)  at (2.3,-0.30)    {$p$};
    \node[redP]             (npq1) at (3.45,-0.30)   {};
    \node[redP,  tokens=1]  (npi2) at (1.15,-1.50)   {};
    \node[T]                (np2)  at (2.3,-1.50)    {$p$};
    \node[redP]             (npq2) at (3.45,-1.50)   {};
    \node[T]                (ns)   at (4.6,0)        {$s$};

    \draw[blueflow] (na)   -- (npo);
    \draw[blueflow] (npo)  -- (ns);
    \draw[redflow]  (na)   -- (npi1);
    \draw[redflow]  (npi1) -- (np1);
    \draw[redflow]  (np1)  -- (npq1);
    \draw[redflow]  (npq1) -- (ns);
    \draw[redflow]  (na)   -- (npi2);
    \draw[redflow]  (npi2) -- (np2);
    \draw[redflow]  (np2)  -- (npq2);
    \draw[redflow]  (npq2) -- (ns);

    \node[ptitle] at (2.3,{\ycapB-0.18}) {(iv) object-centric Petri net 2};
  \end{scope}
  \fi

  \ifnum\mtfbrief=1
    \draw[rule] (5.80,-1.85) -- (5.80,1.65);
    \node[blueP, minimum size=4mm, label distance=2pt, label={[plab]right:order}] at (0.2,\legendY) {};
    \node[redP,  minimum size=4mm, label distance=2pt, label={[plab]right:item}]  at (2.2,\legendY) {};
  \fi

\end{tikzpicture}
\unskip}
  \caption{One process in four notations and its string diagram. Accept an order with two
    items~($a$), pack each item~($p$), and ship once both are packed~($s$), over the object types
    order (blue) and item (red). Panels~(i) and~(iv) are two object-centric Petri nets, whose
    circles are places holding resources. The object-centric causal net of panel~(ii) has no places
    and carries the same resources as markers on its arcs. Circle nodes in the process tree of
    panel~(iii) encode structure rather than resources, and the diamond gateways of the typed BPMN
    model in panel~(v) do the same. Panel~(vi) is the string diagram, which encodes both minimally,
    structure as topology and resources as wires. Each panel is equivalent to panel~(vi),
    reconstructed from the local connectivity of its activities, and all six emit the language
    $\{\langle a,p,p,s\rangle\}$.}
  \label{fig:oc-two-faces}
\end{figure*}

Organisations audit, redesign, and argue over a picture of the process that a process-mining tool
drew from an event log. A second tool given the same log draws a different picture. The pictures
come in several incompatible kinds, and a
modeller given two of them cannot decide whether they describe the same process without first
choosing a translation from one kind into the other. The comparisons that survive such a translation
run over sequences, either the activity sequences a model allows or the states a model passes
through as it produces them, and both readings replace concurrent activities by the orders in which
they might have occurred. Here we give the four notations used in practice a single common form in
which the concurrency survives, and a test that decides whether two models, drawn in the same
notation or in two, have the same structure. The test computes for each model a signature, a
finite description of which activities hand which objects to which activities, and compares the
two signatures for equality.

No comparison in current use decides whether two models describe the same process. Comparing
trace languages conflates two activities that run concurrently with an exclusive choice between
their two orders, because both allow the same two sequences. The transition systems behind those
sequences make the same identification, one interleaved state graph for the concurrent pair and
one for the choice, and the two graphs are
bisimilar~\cite{sassoneModelsConcurrencyClassification1996}, so bisimulation inherits the
conflation. The equivalences that do see the difference, among them history-preserving
bisimulation~\cite{vanglabbeekRefinementActionsEquivalence2001}, compare runs as partial orders,
but each is defined inside a single notation's unfolding semantics and settled by a matching game
over its state space, so applying one across notations requires the very translation whose
correctness is in question. Replay-based conformance checking measures a model against a log
rather than against a second model. None of the four routes settles the question.

The reason is shared by all four. Every notation records a process through routing devices of its
own, the constructs that steer objects between activities. We take the four that
process-discovery tools return, and between them they exhaust the devices in use, namely a buffer drawn
as a node, a buffer folded into an edge, an operator node, and a gateway.

\emph{Object-centric Petri nets}~\cite{vanderaalstDiscoveringObjectCentricPetri2020} hold objects
in typed places and count them on arc weights. We assume a net is type-consistent, so that no chain
of silent transitions joins places of different object type
(Definition~\ref{def:type-consistent}), which nets produced by object-centric discovery satisfy
automatically. \emph{Object-centric causal nets}~\cite{lissObjectCentricCausalNets2025} have no
places at all and record the same fact in binding markers on their arcs. They are the reference
instance of the construction of Section~\ref{sec:general-framework}, the one in which the routing devices appear in their most
collapsed form. \emph{Object-centric process
trees}~\cite{leemansDiscoveringBlockStructuredProcess2013} route through operator nodes, and are
block-structured and hence series-parallel at the level of
occurrences~(Theorem~\ref{thm:pt-as-sp-language}), the assumption that guarantees their soundness and
excludes the non-series-parallel behaviours. \emph{Typed models in the Business Process Model and
Notation} (BPMN)~\cite{BusinessProcessModel} route through gateways. We take a model to be finite
and expand its OR gateways into combinations of AND and XOR, and require nothing else of it.

No routing device is
determined by the behaviour it encodes, so two correct models of one process may use different
devices, and comparing drawings then compares devices. Section~\ref{sec:related-work} takes up
this literature in detail.

Figure~\ref{fig:oc-two-faces} shows the routing devices at their smallest, on the process its
caption describes, drawn in panels (i)--(v) in the four notations a process-mining tool is most
likely to hand back. Write $a$ for accept order, $p$ for pack item,
and $s$ for ship order, over the two object types $\mathrm{order}$ (blue) and
$\mathrm{item}$ (red).

The five drawings agree on the activities and on little else. The object-centric Petri net of
Figure~\ref{fig:oc-two-faces}(i) records the two items in a typed place carrying the arc weight
$2$. Figure~\ref{fig:oc-two-faces}(iv) draws the same process with one transition per item and
every arc weight $1$, so the device of panel~(i) is a property of the drawing alone.
The object-centric causal net of Figure~\ref{fig:oc-two-faces}(ii) has no places anywhere and
records the same fact on its arcs, so the buffer one notation draws as a node the other folds into
an edge. No picture is a special case of another. The process tree of
Figure~\ref{fig:oc-two-faces}(iii) routes on the arity of a parallel operator, and the BPMN model
of Figure~\ref{fig:oc-two-faces}(v) on a pair of gateways, devices of a third and a fourth kind.

Figure~\ref{fig:oc-two-faces}(vi) is the same process with the routing devices removed, and it is
the picture this paper builds on. A box is an activity occurrence and a wire is a typed object
passing from the activity that produces it to the one that consumes it
(Table~\ref{tab:intuition}). Here $a$ sends one $\mathrm{order}$ wire and two $\mathrm{item}$ wires outward, each
$\mathrm{item}$ wire passes through its own $p$ box, and all three arrive at $s$. Nothing else is
drawn, because the picture records only which activity hands which object to which other activity.

A picture of this kind is a string diagram. It is not a fifth notation to model in.
Nothing in it is left to choose, since every routing device is gone and only the typed hand-overs
remain. The string diagram is the normal
form the four notations share, the way regular expressions are compared through their minimal
automata, and it comes with an algebra that says how two diagrams combine.

The reading of a run as a partially ordered set of occurrences is itself old. Partial orders have
modelled concurrency since the sixties~\cite{prattModelingConcurrencyPartial1986}, net theory
reads each run of a Petri net as such an order, and the equivalences above compare those orders
configuration by configuration. The missing step is from semantics to syntax. In net theory the
partial order is something a model means, computed from the model by unfolding it. Here the partial order is
something a modeller writes, the drawing shared by four notations, and equality of drawings is
decided by comparing finite sets rather than by playing a game.

The signature turns this into the promised test. Every model in the four notations reduces to its
signature, the finite record of its activities and their typed inputs and outputs, and the
reduction removes the routing devices and nothing else. Equal signatures certify the same
structure, the same activities with the same typed inputs and outputs, and hence the same trace language,
whichever notations the two models were drawn in, while unequal signatures witness a structural
difference that no drawing convention removes. Included signatures certify included trace
languages in the same way. The comparison fixes no modelling convention, and
it runs on models as discovery tools emit them. Section~\ref{sec:validation} runs it on the output
of two discovery algorithms from one log. The discovered causal net has exactly the ground-truth
signature, and the discovered Petri net's signature strictly contains it, so the comparison names
each behaviour the Petri net adds. The step from structure to behaviour runs one way, since two models of different
structure can still allow the same traces, and the contribution list below states the direction.

Three properties make the string diagram usable as a common form, and
Table~\ref{tab:intuition} records the vocabulary the rest of the paper uses, with a plain reading
for each term.

\begin{description}
  \item[Connectivity, not order.] The wires record which activities exchange which objects, not
    when anything fires. Two activities with no wire between them are concurrent, and no ordering
    is implied or stored.
  \item[Typed resources with multiplicity.] Each wire carries an object type, and two boxes join
    only where their types agree. An order and a pair of items are therefore different data rather
    than differently labelled data, and object-centric coherence follows from the joining rule with
    no further bookkeeping.
  \item[Deformation invariance.] Only the connections matter. The picture may be stretched, and its
    boxes slid past one another, without changing the process it denotes.
\end{description}

Two rules assemble diagrams from smaller ones. Sequential composition $(;)$ joins two diagrams end
to end by matching outgoing wires to incoming wires of the same type. Parallel composition
$(\otimes)$ sets two diagrams side by side with no interaction between them. A box may have several
output wires of one type, which is how $a$ hands its two $\mathrm{item}$ wires to the two packs of
Figure~\ref{fig:oc-two-faces}(vi). Wires may also split and merge, which copies one object to
several activities and joins the copies again.

\begin{table*}[t]
  \centering
  \caption{The vocabulary of the construction, in the order the paper introduces it. Each term is
    given its plain reading here and its formal definition at the section listed.}
  \label{tab:intuition}
  \begin{tabular}{l p{0.52\linewidth} l}
    \hline
    term & plain reading & defined in \\
    \hline
    string diagram & a picture of one process, with a box for each activity and a wire for each
      object passing between activities & Section~\ref{sec:hypergraphs} \\
    wire & one object of one type, running from the activity that produces it to the activity that
      consumes it & Section~\ref{sec:hypergraphs} \\
    cospan & a fragment with a designated input side and output side, so that two fragments
      glue by matching the wires on the sides they meet & Section~\ref{sec:hypergraphs} \\
    pushout & the gluing itself, which identifies the wires two fragments hand to each other and
      keeps everything else & Section~\ref{sec:hypergraphs} \\
    decorated cospan & a cospan carrying data of its own, which the gluing carries along
      & Section~\ref{sec:hypergraphs} \\
    object type & the kind of thing a wire carries, so that two boxes join only where their types agree
      & Section~\ref{sec:hypergraphs} \\
    generator $g_{a,c}$ & one box, meaning an activity together with the types of the objects it
      consumes and produces. An activity gets one box for each context, one way of joining what it
      receives to what it emits &
      Section~\ref{sec:general-framework} \\
    signature $\Sigma$ & the set of a model's boxes, unconnected, each distinguished by its label, its typed
      boundary of incoming and outgoing wires and its constraints on their counts & Section~\ref{sec:hypergraphs} \\
    mediator & a notation's routing device, meaning a place, a gateway, an operator node, or a
      binding set & Section~\ref{sec:general-framework} \\
    $(;)$ and $(\otimes)$ & join two diagrams end to end, and set two diagrams side by side &
      Section~\ref{sec:hypergraphs} \\
    hypergraph category & the setting in which the two joins and wire splitting and merging are
      always available, subject to the axioms of Appendix~\ref{sec:axioms} &
      Section~\ref{sec:hypergraphs} \\
    trace semantics $\Gamma$ & the map sending a diagram to the sequences its partial order admits
      & Section~\ref{sec:hypergraphs} \\
    \hline
  \end{tabular}
\end{table*}

Because routing devices are never recorded, models that differ only in their devices become the
same string diagram. Figure~\ref{fig:routing-collapse} shows the effect within a single notation.
Two object-centric Petri nets returned by two discovery algorithms from one event log, agreeing on
their activity labels and differing in their places and silent transitions, both reduce to the one
string diagram of Figure~\ref{fig:routing-collapse}(iii) (Remark~\ref{rem:ab-place-vs-tau}).

The same picture keeps the structure a sequence comparison loses. The parallel composite
$a\otimes b$, in which two activities are both required with no constraint on their order, and the
exclusive choice between the sequential runs $a\,;\,b$ and $b\,;\,a$ admit the same two sequences
$\langle a,b\rangle$ and $\langle b,a\rangle$ and have bisimilar transition systems. As string diagrams the two are plainly different, one diagram
whose causal order is discrete against two diagrams that are each a total order. For a modeller
the two models are not interchangeable, and no sequence-based comparison warns of it. The reading
extends to choice in general, which Section~\ref{sec:general-framework} makes formal. A
model with choice denotes a family of diagrams, one for each resolution of its choices.

\begin{figure*}[!t]
  \centering
  \resizebox{\linewidth}{!}{\begin{tikzpicture}[petriBase, inline, scale=1.05, every node/.style={transform shape},
    silentT/.style={transition, fill=black!22},
    place/.append style={minimum size=6mm},
    transition/.append style={minimum size=6mm, font=\small},
    plab/.style={font=\footnotesize, black!60, inner sep=1pt},
    ptitle/.style={font=\small\bfseries, inner sep=2pt},
    note/.style={font=\scriptsize, black!55, inner sep=1pt},
    rule/.style={black!30, line width=0.5pt}]

  \def\ytop{1.0}
  \def\ybot{-1.0}
  \def\ycap{-3.05}

  \begin{scope}[shift={(0,0)}]
    \node[blueP, tokens=1] (Ap0) at (0,0)          {};
    \node[T]               (Aa)  at (1.0,0)        {$a$};
    \node[redP]            (Api) at (2.0,\ytop)    {};
    \node[blueP]           (Apo) at (2.0,\ybot)    {};
    \node[T]               (Ab)  at (3.0,\ytop)    {$b$};
    \node[T]               (Ac)  at (3.0,\ybot)    {$c$};

    \draw[blueflow] (Ap0) -- (Aa);
    \draw[redflow]  (Aa)  -- (Api);
    \draw[blueflow] (Aa)  -- (Apo);
    \draw[redflow]  (Api) -- (Ab);
    \draw[blueflow] (Apo) -- (Ac);

    \node[ptitle] at (1.5,\ycap) {(i) OCPN, miner 1};
  \end{scope}

  \begin{scope}[shift={(4.5,0)}]
    \node[blueP, tokens=1] (Bp0)  at (0,0)         {};
    \node[T]               (Ba)   at (1.0,0)       {$a$};
    \node[redP]            (Bpi1) at (1.95,\ytop)  {};
    \node[silentT]         (Bt)   at (2.85,\ytop)  {$\tau$};
    \node[redP]            (Bpi2) at (3.75,\ytop)  {};
    \node[blueP]           (Bpo)  at (2.4,\ybot)   {};
    \node[blueP]           (Bpo2) at (2.85,-1.95)  {};
    \node[T]               (Bb)   at (4.65,\ytop)  {$b$};
    \node[T]               (Bc)   at (4.65,\ybot)  {$c$};

    \draw[blueflow] (Bp0)  -- (Ba);
    \draw[redflow]  (Ba)   -- (Bpi1);
    \draw[redflow]  (Bpi1) -- (Bt);
    \draw[redflow]  (Bt)   -- (Bpi2);
    \draw[redflow]  (Bpi2) -- (Bb);
    \draw[blueflow] (Ba)   -- (Bpo);
    \draw[blueflow] (Bpo)  -- (Bc);
    \draw[blueflow] (Ba.south) to[out=-90,in=180] (Bpo2);
    \draw[blueflow] (Bpo2) to[out=0,in=-90] (Bc.south);

    \node[note, anchor=south] at (2.85,\ytop+0.42) {silent split};
    \node[note, anchor=north] at (2.85,-2.32)      {implicit place};

    \node[ptitle] at (2.3,\ycap) {(ii) OCPN, miner 2};
  \end{scope}

  \draw[rule] (10.35,-2.75) -- (10.35,1.75);

  \begin{scope}[shift={(11.55,0)}]
    \node[morphismblue] (Da) at (0,0)       {$a$};
    \node[morphismblue] (Db) at (2.0,\ytop) {$b$};
    \node[morphismblue] (Dc) at (2.0,\ybot) {$c$};

    \draw[bluewire] (Da.west) -- ++(-0.55,0);
    \draw[redwire]  ([yshift=1.6mm]Da.east) to[out=0,in=180] (Db.west);
    \draw[bluewire] ([yshift=-1.6mm]Da.east) to[out=0,in=180] (Dc.west);
    \draw[redwire]  (Db.east) -- ++(0.55,0);
    \draw[bluewire] (Dc.east) -- ++(0.55,0);

    \node[ptitle] at (1.0,\ycap) {(iii) one string diagram};
  \end{scope}

  \node[blueP, minimum size=4mm, label distance=2pt, label={[plab]right:order}] at (0.2,-3.95) {};
  \node[redP,  minimum size=4mm, label distance=2pt, label={[plab]right:item}]  at (2.2,-3.95) {};

\end{tikzpicture}
\unskip}
  \caption{One behaviour over two object types, receive order~($a$), pick item~($b$), and
    invoice~($c$). Panels (i) and (ii) are object-centric Petri nets as two discovery algorithms
    emit them, differing only in a silent transition splitting the item place and an implicit
    place duplicating the order constraint. Both reduce to the single string diagram (iii).}
  \label{fig:routing-collapse}
\end{figure*}

The rest of the paper builds the mathematics and the machinery this picture needs. We work in a
hypergraph category, the setting of compositional category
theory~\cite{fongHypergraphCategories2018} in which wires may split and merge, and we assign to each
notation a signature $\Sigma$, the collection of boxes its models are built from. Places,
silent transitions, gateways and tree operators are all mediators, and contribute nothing to it. Every notation is then presented in
one vocabulary and joined by one composition algebra, so comparing and translating models across
notations becomes well defined.

The paper makes three lines of contribution, each stated as a modelling deliverable.

\begin{itemize}

  \item \textbf{Four canonical presentations as one category.} Every model in each of the four
    notations reduces to a signature, and one construction~(Section~\ref{sec:general-framework})
    does it in all four cases. Formally,
    object-centric Petri nets~(Theorem~\ref{thm:petri-to-cospan-signature}), object-centric causal
    nets~(Theorem~\ref{thm:causal-to-cospan-signature}), object-centric process
    trees~(Theorem~\ref{thm:ptree-to-cospan-signature}), and typed
    BPMN~(Theorem~\ref{thm:bpmn-cospan}) each receive a canonical presentation as a cospan-algebra
    signature. BPMN needs no structural precondition, and its OR
    gateways are handled by finite expansion into AND and XOR combinations. The presentation is
    notation-independent, so the same behaviour receives the same signature whichever notation
    expresses it (Remark~\ref{rem:ptree-recovers-pn}).

  \item \textbf{An equivalence certificate more detailed than traces.} Two models can be checked for
    agreement by comparing their signatures, with no modelling convention to fix first. Equal
    signatures certify the same structure~(Lemma~\ref{lem:sig-complete}) and hence the same trace
    language~(Theorem~\ref{thm:signature-equivalence}), and the certificate separates the genuine
    concurrency $a\otimes b$ from the exclusive choice that trace comparison and bisimulation both
    conflate with it. Included signatures certify included trace
    languages~(Corollary~\ref{cor:signature-inclusion}).
    Because the maps preserve $(;)$ and $(\otimes)$, the verdict persists when a fragment is
    composed into a larger model. The implication runs in one direction only and we do not claim its
    converse.

  \item \textbf{Object-centric by construction.} Object types live on the wires. Generator-cospan
    boundaries carry object-type labels, and type-matching composability enforces multi-object
    coherence structurally, with no reachability bookkeeping, so a model that mixes types cannot be
    assembled wrongly in the first place. The classical untyped notations are the one-type
    specialisation.

\end{itemize}

As a byproduct we prove that the sequence and parallel fragment of process trees generates the
finite series-parallel posets, and only those, at the level of
occurrences~(Theorem~\ref{thm:pt-as-sp-language}, Corollary~\ref{cor:pt-sp-iff}).

The construction covers notations that are procedural and local, in which a model is assembled from
pieces each wired to its immediate neighbours. Globally-constrained declarative languages such as
Declare~\cite{diciccioGeneratingEventLogs2015} or Dynamic Condition Response (DCR)
graphs~\cite{deboisDeclarativeProcessMining2017}, whose constraints range over whole traces, fall
outside it.

Section~\ref{sec:hypergraphs} develops the
categorical background and the trace semantics $\Gamma$ that connects string diagrams to process
behaviour, and Section~\ref{sec:general-framework} presents the single construction the four
notation sections instantiate. The construction theorems appear in Sections~\ref{sec:petri-nets},
\ref{sec:causal-nets}, \ref{sec:process_trees}, and~\ref{sec:bpmn}, each exercised on a small
object-centric model that reduces to its cospan signature. Section~\ref{sec:validation} reports
the recovery study, in which each of two discovered models is compared against a known
ground-truth signature. The full pipeline, from log to signature to verdict, is implemented in the
\texttt{proc-posets} package released alongside the paper. Section~\ref{sec:conversion-framework} develops cross-notation equivalence and
conversion, Section~\ref{sec:related-work} places the work, and Section~\ref{sec:conclusion}
concludes. Appendix~\ref{sec:axioms} collects the hypergraph-category axioms, and
Appendix~\ref{sec:pt-extras} develops the series-parallel representation theorem for process
trees.
